\addchap{Abstract}

Axions and axionlike particles (ALPs) are theoretically well-motivated candidates for the dark matter that constitutes approximately 27\,\% of the total energy content of the universe.
Through the nuclear gradient coupling, an oscillating ALP dark-matter field acts as an effective pseudomagnetic field on nuclear spins, driving resonant spin precession at the ALP Compton frequency.
Nuclear magnetic resonance (NMR) is among the most sensitive techniques for detecting weak oscillating fields, making it a natural tool for direct ALP dark-matter searches.

This thesis presents a theoretical and experimental investigation of ALP dark-matter searches using precision NMR in the context of the Cosmic Axion Spin Precession Experiment (CASPEr).
On the theoretical side, the signal amplitude and optimal measurement strategy for a frequency-scanning search are derived analytically.
Depending on the relative magnitudes of the NMR coherence time $T_2^*$, the intrinsic transverse relaxation time $T_2$, and the ALP coherence time $\tau_a$, four distinct signal regimes are identified, each with a characteristic sensitivity and optimal dwell-time allocation.
In all regimes, the sensitivity to the ALP-nucleon coupling strength scales as $T_\mathrm{tot}^{-1/4}$ with total measurement time, motivating improvements in spin polarization and relaxation times over simply accumulating more data.
The optimal frequency step is equal to the sensitive bandwidth of a single measurement, and the dwell time should be distributed uniformly across frequency steps in regimes~(1)--(3), or proportional to $\nu^{-1}$ in regime~(4).

On the experimental side, the scanning framework is applied in a proof-of-concept search for the ALP-proton gradient coupling using the CASPEr-Gradient-Low-Field apparatus at the Johannes Gutenberg University Mainz.
A thermally polarized liquid methanol sample is used to scan a 240\,Hz band centered at 1.348570\,MHz (proton Larmor frequency), corresponding to ALP masses in the range $5.576741$--$5.577733\,\mathrm{neV}/c^2$.
Three sequential scans were performed; no ALP signal candidates survived a cross-scan consistency check.
A 95\,\% confidence-level upper limit of $|g_\mathrm{ap}| \lesssim 3\times10^{-2}\,\mathrm{GeV}^{-1}$ is set on the ALP-proton gradient coupling in this mass range, constituting the first dedicated laboratory search therein.

The sensitivity of the current measurement is dominated by the small thermal polarization of the methanol sample ($\approx 1.8\times10^{-7}$).
Projected sensitivity estimates demonstrate that hyperpolarized samples---parahydrogen-induced polarization, dynamic nuclear polarization, or spin-exchange optical pumping of liquid ${}^{129}$Xe---can improve the coupling sensitivity by several orders of magnitude, opening a path toward new physics beyond current bounds.
